\subsection{Benchmark selection and comparison contracts}
The Lehigh University Benchmark (LUBM)~\cite{guo2005} is the first external workload because its maintainers supply
an ontology, deterministic generator, queries, and extensional reference
answers~\cite{lubmwebsite}. UOBM~\cite{ma2006} increases expressive and structural
demands; OWL2Bench~\cite{singh2020} provides explicit OWL 2 profile coverage.
ORE~\cite{parsia2017} supplies real-ontology reasoning tasks. WatDiv
\cite{aluc2014} and BSBM~\cite{bizer2009} instead stress RDF query processing.
Table~\ref{tab:benchmarks} identifies the relevant contracts rather than
treating these suites as interchangeable rankings.

For relational workloads, the Transaction Processing Performance Council's
Benchmark H (TPC-H) evaluates analytical queries over business
data~\cite{tpch}. Our use of selected query components is described below.

\begin{table}[ht]
\caption{Established benchmark candidates and their applicability. Selection is
our assessment; benchmark purpose and language come from the cited sources.}
\label{tab:benchmarks}
\small
\begin{tabularx}{\linewidth}{@{}p{.18\linewidth}YY@{}}
\toprule
Benchmark & Appropriate use & Required boundary \\
\midrule
LUBM & Loading, materialization, exact named query answers & Official ontology and generator; DLP L3 for the unmodified ontology here. \\
UOBM & More connected and expressive ontology workloads & Pin the generator: Oxford's version changes distributions; reject unsupported logical features. \\
Oxford maintenance suite & Compare DRed, DRed$^c$, B/F, B/F$^c$ & Match the exact transformed rule program and fact-deletion batches. \\
ORE 2015 & Coverage, consistency, classification, realization & Record unsupported ontologies; a filtered DLP track cannot inherit full DL/EL rankings. \\
OWL2Bench & Feature coverage and scaling by OWL 2 profile & Property chains and some other profile features exceed this compiler's contract. \\
WatDiv / BSBM & Join diversity and RDF query services & Separate query adaptation from inference; not tests of DRed or OWL completeness. \\
\bottomrule
\end{tabularx}
\end{table}

\subsection{Official LUBM(1,0): independent answer validation}
The initial external study uses UBA 1.7's unchanged bundled generator bytecode,
university index zero, seed zero, and one university. Fifteen department files
are generated. The original ontology IRI is retained; imports are resolved
locally and their 15 document import statements removed after resolution.
No logical ontology axiom is projected away. The resulting graph has 100,851
distinct triples, including 293 ontology triples. L2 rejects existential
consequents; DLP L3 accepts the original ontology and completes materialization.

The official query file contains punctuation errors and an older namespace.
The adapter records both the upstream bytes and its normalized queries, fixing
only those syntax/namespace differences. RDFLib parses each basic graph
pattern; the reasoner's existing indexed join evaluator executes it. This is
an explicit BGP adapter, not a complete SPARQL frontend. Bindings include named
individuals and literal values, matching the official extensional answers.
Skolem witnesses remain present in the materialization and in intermediate
reasoning.

All nine fresh workers (three configurations, three balanced blocks) return
exactly the official answer tuples for all fourteen queries: 126 successful
query checks, including both missing and excess tuple detection. In query
order, the cardinalities are
\begin{equation*}
(4,0,6,34,719,7790,67,7790,208,4,224,15,1,5916).
\end{equation*}
The full internal closure contains 234,480 facts; its named Univ-Bench projection
contains 138,478. All nine full internal closure hashes also agree. Internal
and named closure hashes are recorded separately; agreement does not independently certify every
unqueried existential consequence. Reference-answer equality is the stronger
external evidence for the named benchmark questions.

\begin{table}[ht]
\caption{Official LUBM(1,0), three fresh processes per configuration. Loading/reasoning times are seconds; the fully consumed fourteen-query total is milliseconds. Earlier current state is commit 5531be4.}
\label{tab:lubm}
\centering\small
\begin{tabular}{@{}lrrrr@{}}
\toprule
Configuration & Parse + reasoner & Compile & Materialize & Queries (ms) \\
\midrule
b1254c4 & 7.423 & 1.379 & 3.322 & 217.714 \\
current-fixed & 7.218 & 1.362 & 3.180 & 151.097 \\
current-adaptive & 7.387 & 1.395 & 3.310 & 150.822 \\
\bottomrule
\end{tabular}
\end{table}

Here ``current'' denotes the earlier \code{5531be4} research state, before the
subsequent optimizations motivated by the extended literature review. It is
preserved as a separate external control. The fixed variant has paired median
ratios of 0.957 for materialization and 0.6925 for the sum of fourteen query
times against \code{b1254c4}. This is a local same-host comparison, not a
comparison with the thesis's historical software. Three blocks and one
university establish neither robust scalability nor an external performance
frontier.

Query preparation is outside the query timer, and each query is fully consumed
once per worker. Cold denotes a fresh Python process, not a flushed operating
system filesystem cache. Compilation and materialization are subintervals of
reasoner construction; they must not be added to the already inclusive
parse-plus-reasoner column. Peak RSS includes RDF export and validation and is
approximately 984 MiB. The certificate mechanism concerns retractions and is
inactive in this static benchmark.

\subsection{Published reference results and what they establish}
The 2005 LUBM study~\cite{guo2005}, Table 1 (author PDF p.~15), reports
LUBM(1,0) loading times of 13 s for Sesame-Memory, 542 s for Sesame-DB, and
343 s for DLDB-OWL. Table 3 (author PDF p.~31) gives Q1 latencies of 15, 46,
and 59 ms respectively, all with four answers. Its 103,397-triple ABox count
includes duplicates. The machine was a 1.8 GHz Pentium 4 with 256 MB RAM,
Windows XP, Java 1.4.1, and a 512 MB heap; Sesame 1.0 used RDFS inference and
DLDB-OWL used its 2004-03-29 release. These loading, storage, and semantic
configurations differ from our DLP L3 materialization and BGP adapter. The numbers
are historical reference points and must not be divided by our times.

For the new maintenance hypothesis, the Oxford workloads are more relevant
than a query-only benchmark. Table~\ref{tab:published-maintenance} records two
verified rows from Hu et al.~\cite{hu2018}, Section 6 and Table 1 (author PDF
pp.~7--8). That study averages ten randomly selected 1,000-fact deletions on a
Dell PowerEdge R720 with two 2.6 GHz Xeon E5-2670 processors, 256 GB RAM, and
Fedora 24. It demonstrates why plain DRed alone is an insufficient competitive
control for a certificate optimization.

\begin{table}[ht]
\caption{Published small-deletion reference times from Hu et al.~\cite{hu2018}.
Seconds, on their hardware and transformed programs; these are not new runs.}
\label{tab:published-maintenance}
\centering\small
\begin{tabular}{@{}lrrrrr@{}}
\toprule
Program & Explicit facts & DRed & DRed$^c$ & B/F & B/F$^c$ \\
\midrule
Reactome-U & 12.5 million & 1.03 & 0.06 & 1.00 & 0.05 \\
UOBM-U & 254.8 million & 1,185.87 & 179.24 & 31.96 & 1.18 \\
\bottomrule
\end{tabular}
\end{table}

The U transformation is complete but unsound relative to the original OWL
ontology; R is sound but incomplete, and Claros-LE adds manually chosen rules.
These labels describe OWL approximation, not an error in Datalog evaluation of
the transformed program. The original counting artifact server could not be
retrieved during this study. No version-identical rerun of those Oxford data
or external DRed$^c$/B/F$^c$ executable is claimed.

Other widely used benchmarks expose different limitations. ORE evaluates
consistency, classification, and realization under DL/EL task limits; filtering
to supported DLP inputs changes the competition population. OWL2Bench includes
features such as property chains that this compiler currently rejects. WatDiv
uses diverse SPARQL patterns and BSBM includes service-level and aggregate query
workloads. Successful Horn closure does not implement those complete interfaces.
Their proper role is a separate coverage or query track, with unsupported
features and conversion costs recorded.

The final published RPT+ evaluation by Qiao, Boncz, and Zhang~\cite{qiao2026}
uses DuckDB 1.3.0 and eight threads on two Xeon Platinum 8474C processors with
512 GB DDR5 memory under Debian 12. Its benchmark population includes 18
TPC-H join queries at SF100, 113 Join Order Benchmark queries, eight Appian
queries, and 18,251 SQLStorm queries (13,308 completed, 2,837 failed, and 2,106
reached a ten-second timeout). Reported geometric-mean speedups are 1.17 for
TPC-H, 1.47 for JOB, 1.01 for Appian, and 1.28 for SQLStorm. The final PVLDB
paper is the reference here; an earlier author-hosted draft has different
numbers. Our smaller positive join components do not measure those full SQL
workloads, and their local speed ratios cannot be ranked against these values.

The pinned ZodiacEdge repository~\cite{zodiacartifact} reports LUBM(1) initial,
rule-insertion, and rule-deletion times of 588.9, 216.9, and 252.5 ms for the
$*$ subset, and 603.7, 56.3, and 21.5 ms for the $\#$ subset. These are author
README observations, separate from the paper's wind-turbine workloads with
aggregation and negation~\cite{xu2023}. Hardware, generation seed, repetition
count, and exact timing boundaries are not supplied for these README rows.
The supplied positive LUBM program nevertheless permits a stronger test: run
the author's unmodified engine and our implementation on one host with the
same facts, rules, update sequence, and validated semantic outputs.
